\begin{abstract}\phantomsection\label{sec:resumen}
\noindent Let \(Q(n)\) be the number of ways to place \(n\) mutually nonattacking queens on an \(n\times
n\) board. The canonical predictor \(Q_{\rm pred}(n)\), a finite
formula built from the maximizer of an entropy on the board, reproduces the leading term and the constant
term of \(\log Q(n)\); both facts come from earlier preprints by the author. We now determine the first
correction to the finite predictor: the signed error
\(R_Q(n)=\log\bigl(Q(n)/Q_{\rm pred}(n)\bigr)\) satisfies \(nR_Q(n)\to\cerr\). We give an explicit
formula for this limit and obtain \(Q(n)/Q_{\rm pred}(n)=1+\cerr/n+o(1/n)\).

\smallskip\noindent To obtain the formula, we choose a solution uniformly at random and reveal its rows
in random order. The limit combines an initial correction, the integral of the mean defect left by each
insertion, and three contributions from the board that remains at the stopping time: the rook
contribution, that of a pair of board squares on the same diagonal, and that of the interaction between
the two
diagonals of a square. Comparing two stopping levels gives the limit and proves that the integral
converges absolutely. The proof relies on results from those preprints and is computer-assisted at two
steps.

\smallskip\noindent Evaluating the formula numerically gives \(\cerr\approx0.40550038\), with an empirical
error estimate of \(3\cdot10^{-8}\) and a conservative band \(0.4055004\pm2\cdot10^{-7}\). These are
numerical estimates, not certified enclosures: the sign and decimal digits of \(\cerr\) remain to be
certified.

\medskip\noindent
\textbf{Keywords:} $n$-queens problem; asymptotic enumeration; permanents; rook problem;
interval arithmetic.

\noindent
\textbf{MSC 2020:} 05A16; 60C05; 15A15.
\end{abstract}
